\section{Introduction}
\label{sec:intro}

\begin{figure}[t]
    \centering
    \includegraphics[width=\linewidth]{fig/fig1.pdf}
    \caption{\textbf{The diagram of the proposed GenSplat.} Given expert demonstrations, GenSplat employs a feed-forward 3D reconstruction pipeline to render geometrically consistent novel viewpoints under controlled perturbations, explicitly boosting camera pose diversity. Experiments demonstrate that policies trained on GenSplat-augmented data achieve competitive performance across different perturbation levels, substantially outperforming those trained solely on human-collected demonstrations.
    % Given a sequence of demonstration trajectories from a fixed camera, GenSplat reconstructs each frame with feed-forward 3D Gaussian splatting and renders geometrically consistent novel viewpoints in under 0.13s. The original and augmented demonstrations are combined to train visuomotor policies, which benefit from increased viewpoint diversity and achieve robust generalization to unseen camera poses.
    % Given human demonstrations, GenSplat employs a feed-forward reconstruction method to render novel viewpoints, generating a comprehensive dataset for policy learning. The synthesized views are augmented by varying perturbation levels, enhancing the policy's robustness to changes in camera viewpoints. Experiments in real-world show that policies trained on synthetic data can achieve competitive performance across different perturbation levels, outperforming policies trained solely on human-collected data.
    }
    \vspace{-2ex}
    \label{fig:motivation}
\end{figure}

% Developing policies that execute reliably in complex, real-world environments has long been a central pursuit in robotics and Embodied AI.
% Recent advances in large-scale end-to-end imitation learning have shown considerable promise~\cite{zhao2025cot,black2410pi0,zitkovich2023rt,brohan2022rt,team2024octo,driess2023palm}.
% Despite this progress, they hinge on expansive, diverse datasets~\cite{o2024open,khazatsky2024droid,walke2023bridgedata,dasari2019robonet} and still require time-consuming, manual data collection when deploying to novel environments or robotic systems.
% This reliance highlights a critical vulnerability in the current observation-to-action paradigm:
% Even slight shifts in camera pose can markedly degrade performance by altering occlusions, apparent scale, and the perception of action-relevant visual cues~\cite{pumacay2024colosseum,gao2024efficient,abouzeid2025geoaware,pang2025learning}.
% A common response is to re-collect demonstrations to obtain fresh observation–action sequences for every novel viewpoint, which drastically multiplies teleoperation time and coordination costs~\cite{sadeghi2018sim2real}.
% This dilemma naturally raises the question: \textbf{How can we efficiently augment camera pose in real-world deployments without additional data collection?}

While recent advances in large-scale imitation learning have demonstrated considerable promise in Embodied AI~\cite{zhao2025cot,black2410pi0,zitkovich2023rt,brohan2022rt,team2024octo,driess2023palm}, their efficacy remains strictly contingent upon expansive, diverse datasets~\cite{o2024open,khazatsky2024droid,walke2023bridgedata,dasari2019robonet}. This data dependency exposes a fundamental vulnerability in the prevailing observation-to-action paradigm: current 2D-centric policies exhibit a pronounced deficiency in novel view generalization. Specifically, because these models inherently rely on perspective-dependent visual representations, even minor extrinsic $SE(3)$ perturbations disrupt the spatial grounding of action-relevant cues~\cite{pumacay2024colosseum,gao2024efficient,abouzeid2025geoaware,pang2025learning}. Mitigating this spatial fragility through the exhaustive re-collection of demonstrations for unseen viewpoints imposes prohibitive marginal costs on teleoperation, thereby hindering scalable real-world deployment~\cite{sadeghi2018sim2real,khazatsky2024droid}. This dilemma naturally motivates a critical inquiry: \textit{How can we augment the observational manifold to induce view generalization without additional human collection?}


To address this, prior efforts have leveraged novel view synthesis (NVS) to artificially expand camera pose diversity, broadly falling into two paradigms: 2D diffusion-based generation~\cite{seo2024genwarp,van2024generative,yang2023movie,zhou2025stable,sargent2024zeronvs} and 3D Gaussian-based reconstruction~\cite{yang2025novel,torne2024robot,li2024robogsim,yu2025real2render2real}.
\textbf{Generative methods} employ diffusion models for plug-and-play viewpoint synthesis, bypassing the need for explicit 3D reconstruction or rigid camera calibration~\cite{yang2023movie,chen2024rovi,tian2024view,shi2025nvspolicy}. Although recent adaptations incorporate 3D-aware conditioning~\cite{sargent2024zeronvs,lu2025movis}, their optimization objectives remain strictly confined to the 2D image plane. Consequently, these generated augmentations lack a unified metric workspace, frequently violating epipolar geometry and multi-view visibility constraints. 
This manifests as severe scale drift, erroneous parallax, and implausible occlusions, which are fatal for precise robotic manipulation. Moreover, underutilizing the multi-camera topology of robotics datasets restricts both geometric fidelity and synthesis controllability.


In contrast, \textbf{Gaussian-based methods} parameterize explicit 3D scenes using anisotropic Gaussian primitives~\cite{kerbl20233d,PapantonakisReduced3dgs}. Rendering from this shared, underlying geometry inherently guarantees strict cross-view physical consistency by explicitly preserving exact occlusion ordering, multi-view parallax, and metric scale, all while enabling real-time synthesis. However, classical 3DGS pipelines suffer from restrictive operational bottlenecks: achieving high-fidelity reconstruction typically necessitates dense multi-view coverage, precise camera pose calibration, and prohibitively expensive per-scene optimization~\cite{yang2025novel,torne2024robot,yu2025real2render2real,duan20244d,li2024robogsim}. Crucially, real-world robotic deployments are fundamentally characterized by extreme viewpoint sparsity, entirely uncalibrated observation setups, and the necessity for instant adaptation across diverse environments. Motivated by this domain gap, we propose a feed-forward reconstruction paradigm capable of instant, high-fidelity scene synthesis from sparse, uncalibrated observations, circumventing the limitations of traditional per-scene optimization.

In this work, we introduce GenSplat, a feed-forward 3D Gaussian Splatting framework designed to expand the observational manifold via novel view synthesis, thereby inducing robust view generalization in visuomotor policies. Circumventing the prohibitively expensive per-scene tuning of classical pipelines, GenSplat leverages recent advances in feed-forward 3DGS~\cite{ye2024no,jiang2025anysplat} to parameterize 3D scenes from sparse, uncalibrated inputs in a single forward pass. This enables efficient, physically grounded multi-view synthesis, systematically expanding the critical dimension of viewpoint diversity required by modern imitation learning architectures~\cite{saxena2025matters,black2410pi0,chi2023diffusion}. However, directly applying naive feed-forward mechanisms to the extreme sparsity of robotic datasets~\cite{o2024open,dasari2019robonet,khazatsky2024droid,walke2023bridgedata} often precipitates severe topological fragmentation and inconsistent occlusions~\cite{kerbl20233d,PapantonakisReduced3dgs}. While such photometric artifacts might be tolerated in pure visual rendering, they severely corrupt the geometric consistency mandatory for reliable visuomotor supervision. To overcome this structural collapse, we propose a 3D-prior distillation objective that extracts geometric knowledge from a 3D foundation model~\cite{wang2025pi} to explicitly regularize the reconstruction landscape. Consequently, GenSplat yields metrically consistent viewpoint augmentations that systematically enhance downstream policy robustness without incurring additional data collection costs.
In summary, our primary contributions are:

(1) We propose GenSplat, a feed-forward 3DGS framework that enables geometrically consistent view augmentation for robotic datasets with sparse and uncalibrated camera observations. The design effectively eliminates the need for sensor calibration, physical viewpoint adjustment, and per-scene optimization.
    
(2) We introduce a 3D-prior distillation strategy that transfers geometric knowledge from a pretrained 3D foundation model to regularize reconstructions, ensuring that synthesized novel views are geometrically consistent and physically reliable for policy learning.

(3) Extensive real-world experiments demonstrate that policies trained with GenSplat-augmented data, achieve strong generalization to unseen camera poses. GenSplat significantly enhances policy performance while demonstrating superior computational efficiency compared to prior methods.

